\documentclass[12pt]{article}
\usepackage{amsmath, amssymb, enumitem, fancyhdr, graphicx, libertine, nth,
  pgfplots, tikz, xcolor}
\pgfplotsset{compat=1.11}
\usepackage[margin=1in]{geometry}
\usepackage[libertine]{newtxmath}
\pagestyle{fancy}
\definecolor{solution-color}{HTML}{000080}
\chead{\emph{EC201 -- Extra Practice Problems}}

% Define new topic command
\newcommand{\topic}[1]{\medskip\begin{center}\large\textsc{#1}\normalsize
  \end{center}}
% Define question counter:
\newcounter{question}[section]
\newcommand{\question}[1]{\refstepcounter{question}
  \subsubsection*{Question~\thequestion~-- #1}}

% Subquestions are enumerate with letters
\newenvironment{subquestions}{
  \begin{enumerate}[label=(\roman*)]}{\end{enumerate}}

\begin{document}
\question{Demand}
Your demand functions for goods 1 and 2 are:
\begin{align*}
  x_1\left( p_1,p_2,m\right)&=
  \begin{cases}
    \frac{2m - 100}{p_1} &\text{ if } m \geq 50 \\
    0 & \text{ if } m < 50 \\
  \end{cases} \\
  x_2\left( p_1,p_2,m\right)&=
  \begin{cases}
    \frac{100 - m}{p_2} & \text{ if } m \leq 100 \\
    0 & \text{ if } m > 100 \\
  \end{cases}
\end{align*}
Assume that your income is always such that $50\leq m\leq 100$ so we can write demand as:
\begin{align*}
  x_1\left( p_1,p_2,m\right)&=\frac{2m - 100}{p_1} \\
  x_2\left( p_1,p_2,m\right)&=\frac{100 - m}{p_2}
\end{align*}
\begin{subquestions}
  \item Is good 1 an ordinary or a Giffen good?
  \item Is good 2 an ordinary or a Giffen good?
  \item Is good 1 a normal good or an inferior good?
  \item Is good 2 a normal good or an inferior good?
  \item Is good 1 a substitute, a complement, or neither, for good 2?
\end{subquestions}

\question{Production}

A firm cuts down enormous trees and sells the wood. To cut down a tree it needs 2
sawyers (workers who saw down trees for a living) and a giant two-person saw.
One sawyer is not able to use a saw on their own, as the saws are very big.  If
there are 3 sawyers with only 1 saw, the third sawyer will be idle (they
can't help as the saw can only be used be two people).

\begin{subquestions}
  \item Call the sawyers $x_1$ and the giant saws $x_2$.
    Given the information provided, what is the production function $f\left(
    x_1,x_2 \right)$ for cutting down trees?
  \item Sketch an isoquant for this production function.
\end{subquestions}


\question{Uncertainty and Monopoly}
You are considering developing a new piece of software currently unavailable in
the world. If you develop it, you will patent it and then you will have monopoly
rights to sell that software.
You have \$200 sitting in your bank account and the development of the software
will cost \$100. Your friend, who is a marketing expert, tells you that the
demand for your software will be:
\[
  D\left( p \right)=50-\frac{1}{2}p,
\]
where $p$ is the
price per download of the software. Each additional unit of software sold
will incur you no cost.

However, you are concerned that the government will regulate the sale of your
software.  Your friend in the government gives you the following information.
\begin{itemize}
  \item There is a $\frac{1}{3}$ probability that the government will not
    regulate you at all, and you will be free to choose any price you wish.
  \item There is a $\frac{1}{3}$ probability that the government will force you
    to charge at marginal cost.
  \item There is a $\frac{1}{3}$ probability that the government will force you
    to charge a price such that you will break even.
\end{itemize}
Answer the following questions:
\begin{subquestions}
\item What profit will you make from each of the three
  possible government actions (no regulation and two types of regulation)?
  \item If your utility for wealth is $u\left( W
    \right)=\sqrt{W}$, will you decide to go ahead and develop the software?
    Show why or why not.
\end{subquestions}

\question{Price Discrimination}
An art gallery is frequented by both art students and rich art lovers.
The inverse demand curve for both groups for entry into the art gallery is:
\begin{align*}
  p_S\left( q_S \right) &= 20-q_S \\
  p_R\left( q_R \right) &= 40-q_R
\end{align*}
The art gallery has the following cost function:
\[
  c\left( q_S,q_R \right)=10\left(q_S + q_R\right) +10
\]
\begin{subquestions}
  \item Suppose first that the art gallery can't tell the
    two groups apart. What price should it charge for entry? How many people
    will visit the art gallery? What will the art gallery's profits be?
    What is the consumer surplus?
  \item Now suppose the art gallery can charge different
    prices to the different groups. The students will show that they are
    students by showing their student cards. What prices should the art gallery
    charge to both groups? How many of each group will visit the art gallery?
    What will the art gallery's profits be? What will the total consumer
    surplus be?
  \item Suppose the art gallery can't tell the groups apart
    and now decided to offer two types of gallery experiences. One gives only
    entry but the other also includes a tour. The prices and quality of two
    different entry tickets are designed such that the students self select
    into the ticket that doesn't include a tour and the rich group self selects
    into the more ticket that does include the tour. What degree of price
    discrimination would this be? (No calculation is required).
\end{subquestions}
\end{document}
